\documentclass[a4paper]{article}
\usepackage[pdftex]{graphicx}
\usepackage[utf8]{inputenc}
\usepackage{enumerate}
\usepackage{icomma}
\usepackage{amssymb}
\usepackage{tikz}
\usepackage{href-ul}
\hypersetup{
	colorlinks=true,
	linkcolor=blue,
	urlcolor=blue}
\usepackage{geometry}
\geometry{a4paper, top=15mm, left=15mm, right=15mm, bottom=15mm,
	headsep=10mm, footskip=12mm}

\begin{document}

\begin{tikzpicture}(10,3)
	\draw[ultra thick](2,0) --(10,0) -- (10,1.3) --(2,1.3) -- (2,0);
	\draw[fill=black](2,0)-- (0,.6) -- (2,.6) -- (2,0);
	\node at (6,.6) {\large Lösungsblatt von \href{https://www.okuyakl.de}{www.okuyakl.de}};
\end{tikzpicture}
\vspace{0.5 cm}
\noindent

\fbox{
	\begin{minipage}{0.96\textwidth}
	\noindent{\bf Aufgabe 1. a)}\\
	Die allgemeine Geradengleichung ist:
	$$y=m\cdot x + t $$
	Hierbei ist die Steigung $m$:
	$$ m ={y_B - y_A \over x_B -x_A} = {3+1 \over 2+3}=0,8$$
	Durch Einsetzen von $A$ und $m$ erhalten wir für den y-Achsenabschnitt $t$:
	$$ -1 = 0,8 \cdot (-3) +t \quad \Rightarrow \quad t = 1,4$$	
	\end{minipage}
}

\noindent
\fbox{
	\begin{minipage}{0.48\textwidth}
	\noindent{\bf Aufgabe 1. b)}\\	
	\includegraphics[width=8 cm]{gisela042}
	\end{minipage}
}
\fbox{
	\begin{minipage}{0.48\textwidth}
		\noindent{\bf Aufgabe 1. c)}\\
		Punkt P:
		$$
		\begin{array}{rcll}
		5 &=& 0,8 \cdot x_P +1,4 \\
		0,8 x_P &=& 3,6 \\
		x_P &=& 4,5 
		\end{array}
		$$
		
	Punkt Q:
	$$
	\begin{array}{rcll}
	y_Q &=& 0,8 \cdot (-5) +1,4 \\
	y_Q &=& -2,6
	\end{array}
	$$	


\noindent{\bf Aufgabe 1. d)}\\
	Punktprobe:
	$$
	\begin{array}{rcll}
	0,7 &=& 0,8 \cdot (-1)+1,4 \\
	0,7 &=& 0,6 & (f)\\
	&\Rightarrow& E \notin g  
	\end{array}
	$$

	Punktprobe:	
	$$
	\begin{array}{rcll}
		17,4 &=& 0,8 \cdot 20 +1,4 \\
	    17,4 &=& 17,4 & (w)\\
		&\Rightarrow& F \in g  
	\end{array}
	$$	
	\end{minipage}
}
\fbox{
	\begin{minipage}{0.48\textwidth}
	\noindent{\bf Aufgabe 1. e)}\\
	Für die Steigung der  orthogonalen Geraden gilt:
	$$m_2 = -{1 \over m_1} = - {1 \over 0,8} = -1,25$$
	Punkt $C$ in $y=mx+t$:
	$$ 2 = -1,25 \cdot 0 + t \quad \Rightarrow \quad t = 2$$
	Funktionsgleichung:
	$$h: \quad y= -1,25x +2 $$	
	\end{minipage}
}
\hfill
\fbox{
	\begin{minipage}{0.48\textwidth}
	\noindent{\bf Aufgabe 1. f)}\\	
	Es gilt $m=-1,25$
	Punkt $D$ in $y=mx+t$:
	$$ 1 = -1,25 \cdot (-1) + t \quad \Rightarrow \quad t = -0,25$$
	Funktionsgleichung:
	$$h: \quad y= -1,25x - 0,25 $$	
	\end{minipage}
}
\fbox{
	\begin{minipage}{0.48\textwidth}
	\noindent{\bf Aufgabe 1. g)}\\	
	$$
	\renewcommand{\arraystretch}{2}
	\begin{array}{rcll}
		0 &=& 0,8 \cdot x + 1,4 \\
		-1,4 &=& 0,8 \cdot x \\
		x &=& -1,75 & N(-1,75|0)
	\end{array}
	$$
		
	\end{minipage}
}
\hfill
\fbox{
	\begin{minipage}{0.48\textwidth}
	\noindent{\bf Aufgabe 1. h)}\\	
	Punkt $Y$ in $y=mx+t$:
	$$ -1 = 0,8 \cdot 0 + t \quad \Rightarrow \quad t = -1$$
	Funktionsgleichung:
	$$n: \quad y= 0,8x -1 $$	
	\end{minipage}
}
\fbox{
	\begin{minipage}{0.48\textwidth}
	\noindent{\bf Aufgabe 2. a)}\\
	\includegraphics[width=8 cm]{gedrei042}
	\end{minipage}
}
\fbox{
	\begin{minipage}{0.48\textwidth}
	\noindent{\bf Aufgabe 2. b)}\\	
	Die Funktionsgleichung lässt sich sofort angeben:
	$$g_2: \quad y=0,4x + 1,5 $$
	\noindent{\bf Aufgabe 2. c)}\\
	$R$ in $y=0,4x+t$:
	$$-5 = 0,4\cdot (-4,5) + t \quad \Rightarrow t = -3,2$$
	$$g_3: \quad y= 0,4x - 3,2 $$
	\noindent{\bf Aufgabe 2. d)}\\
	Wir setzen A in $y=0,4x+4$ ein:
	$$	4,9\bar{3}  = 0,4\cdot {7 \over 3} +4 = 4,9\bar{3} \quad (w)$$
	Punktprobe positiv: $A \in g_1$
	\end{minipage}
}

\noindent
\fbox{
	\begin{minipage}{0.48\textwidth}
	\noindent{\bf Aufgabe 3. a)}\\
	Steigung $m=0$; $t=5$;
	$$g_4: \quad y=5$$
	\noindent{\bf Aufgabe 3. b)}\\
	Steigung $m=-1$; $t=0$;
	$$g_5: \quad y=-x$$
	\noindent{\bf Aufgabe 3. c)}\\
	Steigung m:
	$$ m ={y_B - y_D \over x_B -x_D} = {-2-1 \over 7-(-3)}=-0,3$$
	y-Achsenabschnitt t:
	$$1 = -0,3\cdot(-3) + t \quad \Rightarrow t = 0,1$$
	Funktionsgleichung:
	$$g_5: \quad y=-0,3x+0,1$$
	\end{minipage}
}
\hfill
\fbox{
	\begin{minipage}{0.48\textwidth}
	\noindent{\bf Aufgabe 4. a)}\\
	Schnitt mit y-Achse = t:
	$$S_y(0|2)$$
	Schnitt mit x-Achse = Nullstelle:
	$$
	\renewcommand{\arraystretch}{2}
	\begin{array}{rcll}
	0&=&-{3\over 4} \cdot x +2 &| -2 \\
	-2 &=& -{3\over 4} \cdot x &| \cdot(-{4 \over 3})\\
	{8 \over 3} &=& x &\Rightarrow S_x({8 \over 3}|0)
	\end{array}
	$$
	\noindent{\bf Aufgabe 4. b)}\\
	$H$ in $g_6$:
	$$
	\renewcommand{\arraystretch}{2}
	\begin{array}{rcll}
	-{1 \over 4}&=&-{3\over 4} \cdot x +2 &| -2 \\
	-2,25 &=& -{3\over 4} \cdot x &| \cdot(-{4 \over 3})\\
	3 &=& x &\Rightarrow H(3|-{1 \over 4})
	\end{array}
	$$
	\end{minipage}
}


\begin{center}
	\includegraphics[width=7 cm]{../../viecher/endcomic.pdf}
	
	Hier geht es zurück zum \href{https://www.okuyakl.de/math/m9gerL042/aa042.pdf}{Aufgabenblatt}
\end{center}


\end{document}

